\documentclass[11pt]{article}

\usepackage[margin=1in]{geometry}
\usepackage[T1]{fontenc}
\usepackage[utf8]{inputenc}
\usepackage{lmodern}
\usepackage{graphicx}
\usepackage{booktabs}
\usepackage{longtable}
\usepackage{array}
\usepackage{xcolor}
\usepackage{hyperref}
\usepackage{float}
\usepackage{caption}
\usepackage{enumitem}

\hypersetup{
    colorlinks=true,
    linkcolor=blue!45!black,
    urlcolor=blue!45!black,
    citecolor=blue!45!black
}

\definecolor{headingblue}{HTML}{1F4D78}
\definecolor{mutedgray}{HTML}{555555}

\setlength{\parskip}{0.55em}
\setlength{\parindent}{0pt}
\setlist[itemize]{topsep=0.25em,itemsep=0.25em}
\setlist[enumerate]{topsep=0.25em,itemsep=0.25em}

\newcommand{\feature}[1]{\texttt{#1}}
\newcommand{\metric}[1]{\textbf{#1}}
\newcommand{\figpath}[1]{figures/#1}

\title{\textbf{Employee Turnover Prediction}\\
\large Final Report: Data Preparation, Modeling Results, and Interpretability}
\author{Prepared for advisor review}
\date{June 25, 2026}

\begin{document}
\maketitle

\begin{abstract}
This report summarizes the final employee-turnover prediction pipeline developed from three HR data files. The final design uses file1 and file2 for model development and validation, while file3 is held out as an external test set. The supplied \feature{leave\_ind} field is used as the prediction target. Aziva exit-code metadata was audited for label consistency, but excluded from modeling because it would leak outcome information. The selected model, \feature{XGBoost\_\_learned\_imputation}, achieved validation PR-AUC 0.2052 and external file3 PR-AUC 0.1462. Since file3 prevalence is 4.87\%, this external PR-AUC is approximately 3.00 times the baseline positive rate. The model is therefore useful primarily as a ranking and prioritization tool, not as an autonomous decision system.
\end{abstract}

\tableofcontents
\newpage

\section{Executive Summary}

The goal of this project is to predict employee turnover, represented by the supplied \feature{leave\_ind} target, and to explain the main risk drivers in a way that can support HR analysis at company, department, and individual employee levels. The final pipeline trains on file1 and file2, validates with employee-grouped validation, and evaluates once on file3 as an external test set.

\textbf{Bottom line.} The selected model is \feature{XGBoost\_\_learned\_imputation}. It achieved validation PR-AUC \metric{0.2052}, validation ROC-AUC \metric{0.7615}, file3 PR-AUC \metric{0.1462}, and file3 ROC-AUC \metric{0.7758}. The file3 positive rate is only \metric{4.87\%}; therefore, the file3 PR-AUC is approximately \metric{3.00x} the baseline prevalence. In the top 20\% highest-risk file3 rows, the model captured \metric{51.63\%} of positives at \metric{12.58\%} precision.

Key conclusions:
\begin{itemize}
    \item The model has useful ranking power, especially for prioritizing a subset of employees for review.
    \item Absolute precision remains modest because turnover is a rare event in file3.
    \item Age, tenure, and data-timing variables are strong descriptive predictors.
    \item Actionable or partially actionable signals remain important: compensation, salary trajectory, workload, illness/absence patterns, contract type, and role/manager context.
    \item The main methodological concerns are source shift between file1/file2 and file3, noisy file2 payment imputation, low event prevalence, and small subgroup sizes.
\end{itemize}

\section{Data Description and Study Design}

The project uses three HR Excel data sources. File1 and file2 serve as development data. File3 is reserved as an external test set and is not used for fitting, imputation training, model selection, or threshold selection. The modeling unit is an employee-period row rather than only a single latest employee snapshot. This preserves temporal information, but it also means that interpretation should be framed as ranking employee-period risk rather than predicting one permanent label per employee.

\begin{table}[H]
\centering
\caption{Final train/validation/test split summary.}
\begin{tabular}{lrrrr}
\toprule
Split & Rows & Employees & Positive rows & Positive rate \\
\midrule
Inner train & 10,400 & 3,104 & 831 & 7.99\% \\
Validation & 2,654 & 777 & 207 & 7.80\% \\
File3 external test & 3,140 & 790 & 153 & 4.87\% \\
\bottomrule
\end{tabular}
\end{table}

The lower prevalence in file3 is important. Precision-recall metrics are prevalence-sensitive: a model can retain reasonable ROC-AUC while still producing lower PR-AUC and lower precision when positives are rarer. Therefore, the file3 PR-AUC should be interpreted relative to the 4.87\% baseline.

\section{Data Quality Issues and Corrective Methods}

\subsection{Target definition and leakage controls}

The supplied \feature{leave\_ind} field is used as the target. A specific issue arose around Aziva codes: in file1 and file2, Aziva code 41 corresponds to the resignation concept associated with \feature{leave\_ind = 1}; in file3, Aziva code 42 corresponds to the same conceptual outcome. Thus, Aziva code 41 in files 1 and 2 is equivalent to Aziva code 42 in file3 for the resignation label.

However, Aziva codes were not used as predictors. They are exit/outcome metadata and would leak post-event information into the model. They were used only to audit target consistency.

Leakage controls:
\begin{itemize}
    \item Dropped outcome metadata: \feature{aziva\_kod}, \feature{aziva\_date}, \feature{aziva\_year}, target-like fields, and \feature{leave\_ind}.
    \item Dropped identifiers and bookkeeping fields from model inputs: \feature{source\_employee\_id}, \feature{fictive\_employee}, \feature{calc\_month}, \feature{source}, and \feature{year\_date}.
    \item Kept \feature{calc\_month} only for chronological ordering and within-period permutation analysis; it is not directly modeled.
\end{itemize}

\subsection{Invalid records and cleaning}

\begin{table}[H]
\centering
\caption{Invalid-record cleaning rules.}
\begin{tabular}{lrrrr}
\toprule
Cleaning rule & Total removed & file1 & file2 & file3 \\
\midrule
Negative average illness & 14 & 0 & 9 & 5 \\
Non-positive average payment & 136 & 0 & 136 & 0 \\
Non-positive average workload & 4 & 0 & 3 & 1 \\
Age below 18 & 11 & 0 & 10 & 1 \\
Negative tenure & 0 & 0 & 0 & 0 \\
\bottomrule
\end{tabular}
\end{table}

Overall, 165 of 16,359 rows were removed, leaving 16,194 usable rows. Only six removed rows were positive. Employees above age 75 were intentionally retained because this is feasible in the domain; age below 18 was treated as invalid.

\subsection{Payment missingness and imputation}

Payment missingness was a major issue, especially in file2. The final pipeline uses a fold-safe similar-employee imputer for payment-related variables. The imputer predicts \feature{avg\_Payment}, \feature{Median\_Payment}, \feature{stdevp\_Payment}, and \feature{salary\_change}. It is fitted only on training data inside each validation fold and never on file3. Missingness indicators are added so that the model can learn that the value was originally missing.

\begin{table}[H]
\centering
\caption{Cross-validated payment imputation diagnostics.}
\begin{tabular}{llrrr}
\toprule
Source & Imputed field & MAE & RMSE & Median absolute error \\
\midrule
file1 & \feature{avg\_Payment} & 1,776.4 & 3,809.4 & 523.6 \\
file1 & \feature{Median\_Payment} & 1,844.2 & 3,978.8 & 544.1 \\
file1 & \feature{salary\_change} & 371.4 & 844.0 & 171.4 \\
file2 & \feature{avg\_Payment} & 3,672.0 & 5,123.0 & 2,738.1 \\
file2 & \feature{Median\_Payment} & 3,758.3 & 5,235.8 & 2,810.3 \\
file2 & \feature{salary\_change} & 603.9 & 1,272.1 & 335.1 \\
overall & \feature{avg\_Payment} & 2,364.8 & 4,271.5 & 890.4 \\
overall & \feature{Median\_Payment} & 2,441.0 & 4,429.2 & 900.2 \\
overall & \feature{salary\_change} & 447.1 & 1,018.6 & 215.7 \\
\bottomrule
\end{tabular}
\end{table}

The imputation diagnostics show that file2 payment imputation is materially noisier than file1. This motivated explicit comparison of learned-imputation, no-payment, and native-missing modeling strategies.

\section{Feature Engineering and Preprocessing}

The final preprocessing pipeline combines raw HR variables with time-safe historical summaries. Historical features use only current and prior rows for the same employee. This avoids future leakage while allowing features such as previous payment, previous illness, deltas from previous values, prior means, prior standard deviations, slopes per month, and observed-fraction-to-date variables.

Important preprocessing choices:
\begin{itemize}
    \item Categorical fields are imputed with a missing category and one-hot encoded inside the fitted pipeline.
    \item Numeric fields are median-imputed when required; an XGBoost native-missing variant was also tested.
    \item Payment imputation is inside the model pipeline, so validation folds do not share imputation state.
    \item Protected or unstable descriptive fields such as gender and marital status were excluded from the final default feature list.
    \item \feature{Tafkid} text was excluded because EDA suggested source instability. Related role/organization codes were interpreted cautiously as contextual predictors, not causal levers.
\end{itemize}

\section{Modeling Strategy}

The final model search compared Logistic Regression, Random Forest, and XGBoost under several payment strategies. Class imbalance was handled only through native estimator mechanisms: \feature{class\_weight} for Logistic Regression and Random Forest, and \feature{scale\_pos\_weight} for XGBoost. Synthetic oversampling such as SMOTE was not used in the final pipeline because the main concern appears to be source shift and mixed-type feature instability rather than only minority-class scarcity.

Payment strategies tested:
\begin{itemize}
    \item \textbf{learned\_imputation:} similar-employee payment imputer fitted within each training fold.
    \item \textbf{native\_missing:} XGBoost receives numeric missingness directly where possible.
    \item \textbf{no\_payment:} payment and salary features, including payment histories, are removed.
\end{itemize}

Model selection was based on validation PR-AUC. File3 was used only after model and threshold selection.

\section{Model Results}

Table~\ref{tab:model-results} summarizes the final model comparison. The selected model is not necessarily the candidate with the single best file3 score; it is the model chosen before looking at file3. This preserves file3 as an external test set rather than an informal tuning set.


\begin{table}[H]
\centering
\caption{Validation performance used for model selection.}
\label{tab:model-results}
\small
\begin{tabular}{p{5.0cm}rrrrr}
\toprule
Candidate & AUC & PR-AUC & F1 & Precision & Recall \\
\midrule
\feature{XGBoost\_\_learned\_imputation} & 0.7615 & 0.2052 & 0.2897 & 0.2027 & 0.5072 \\
\feature{Random Forest\_\_no\_payment} & 0.7361 & 0.1949 & 0.2466 & 0.2072 & 0.3043 \\
\feature{XGBoost\_\_native\_missing} & 0.7593 & 0.1942 & 0.2780 & 0.1895 & 0.5217 \\
\feature{XGBoost\_\_no\_payment} & 0.7623 & 0.1925 & 0.2818 & 0.1805 & 0.6425 \\
\feature{Random Forest\_\_learned\_imputation} & 0.7395 & 0.1876 & 0.2594 & 0.1848 & 0.4348 \\
\feature{Logistic Regression\_\_no\_payment} & 0.7566 & 0.1860 & 0.3140 & 0.2387 & 0.4589 \\
\feature{Logistic Regression\_\_learned\_imputation} & 0.7668 & 0.1859 & 0.2990 & 0.2112 & 0.5121 \\
\bottomrule
\end{tabular}
\end{table}

\begin{table}[H]
\centering
\caption{External file3 performance after validation-based model selection.}
\small
\begin{tabular}{p{5.0cm}rrrrrrr}
\toprule
Candidate & AUC & PR-AUC & F1 & Precision & Recall & Recall@20\% & Precision@20\% \\
\midrule
\feature{XGBoost\_\_learned\_imputation} & 0.7758 & 0.1462 & 0.1913 & 0.1147 & 0.5752 & 0.5163 & 0.1258 \\
\feature{Random Forest\_\_no\_payment} & 0.7442 & 0.1373 & 0.1695 & 0.1493 & 0.1961 & 0.4510 & 0.1099 \\
\feature{XGBoost\_\_native\_missing} & 0.7625 & 0.1254 & 0.2015 & 0.1237 & 0.5425 & 0.5033 & 0.1226 \\
\feature{XGBoost\_\_no\_payment} & 0.7484 & 0.1057 & 0.1703 & 0.0969 & 0.7059 & 0.4444 & 0.1083 \\
\feature{Random Forest\_\_learned\_imputation} & 0.7444 & 0.1736 & 0.2017 & 0.1350 & 0.3987 & 0.4444 & 0.1083 \\
\feature{Logistic Regression\_\_no\_payment} & 0.7508 & 0.1109 & 0.1562 & 0.0866 & 0.7974 & 0.4837 & 0.1178 \\
\feature{Logistic Regression\_\_learned\_imputation} & 0.7265 & 0.1047 & 0.1416 & 0.0771 & 0.8627 & 0.4575 & 0.1115 \\
\bottomrule
\end{tabular}
\end{table}

\begin{figure}[H]
\centering
\includegraphics[width=\linewidth]{\figpath{model_performance_comparison.png}}
\caption{Validation and file3 performance comparison across final candidate models.}
\end{figure}

The figure compares candidate models using validation and external-test metrics. The selected XGBoost learned-imputation model was chosen by validation PR-AUC. Random Forest with learned imputation has a stronger file3 PR-AUC, but selecting it after seeing file3 would be test-set selection.

\section{Feature Importance for the Selected Model}

Interpretability was performed on the locked XGBoost learned-imputation model. Multiple importance methods were used because each answers a different question:
\begin{itemize}
    \item Native XGBoost gain describes how the trees split.
    \item SHAP describes average contribution magnitude.
    \item File3 permutation measures external reliance by shuffling feature groups within period.
    \item Drop-column validation checks whether removing a feature hurts validation performance.
\end{itemize}

These methods are predictive explanations, not causal estimates.

\begin{longtable}{rp{4.0cm}p{3.2cm}p{2.8cm}r}
\caption{Top consensus features for the selected model.}\\
\toprule
Rank & Feature & Domain & Actionability & Consensus \\
\midrule
\endfirsthead
\toprule
Rank & Feature & Domain & Actionability & Consensus \\
\midrule
\endhead
1 & \feature{avg\_Payment} & Compensation & Actionable & 0.894 \\
2 & \feature{Yishuv} & Location / commute & Descriptive & 0.885 \\
3 & \feature{salary\_change\_hist\_prev} & Compensation & Actionable & 0.849 \\
4 & \feature{age} & Age / tenure & Descriptive & 0.837 \\
5 & \feature{avg\_Payment\_hist\_slope\_per\_month} & Compensation & Actionable & 0.822 \\
6 & \feature{tenure\_years} & Age / tenure & Descriptive & 0.817 \\
7 & \feature{avg\_illness\_hist\_vs\_prior\_mean} & Absence / wellness & Partially actionable & 0.816 \\
8 & \feature{illness\_cv} & Absence / wellness & Partially actionable & 0.814 \\
9 & \feature{stdevp\_omes} & Workload & Actionable & 0.810 \\
10 & \feature{contract\_type} & Role / manager / organization & Partially actionable & 0.808 \\
11 & \feature{manager\_Code} & Role / manager / organization & Partially actionable & 0.802 \\
12 & \feature{vetek\_months} & Age / tenure & Descriptive & 0.797 \\
13 & \feature{months\_since\_previous\_record} & Data maturity / timing & Descriptive & 0.796 \\
14 & \feature{stdevp\_Payment} & Compensation & Actionable & 0.787 \\
15 & \feature{manager\_Code\_hist\_prev} & Role / manager / organization & Partially actionable & 0.779 \\
\bottomrule
\end{longtable}

\begin{figure}[H]
\centering
\includegraphics[width=0.95\linewidth]{\figpath{consensus_top_features.png}}
\caption{Consensus feature ranking across available importance methods.}
\end{figure}

This graph aggregates native gain, SHAP, permutation, and drop-column evidence where available. It shows payment, age/tenure, salary-change history, location, illness, workload, and role/contract context in the selected model. Because not every method is available for every feature, the consensus score should be interpreted as a ranking aid rather than a formal statistical test.

\begin{figure}[H]
\centering
\includegraphics[width=0.95\linewidth]{\figpath{shap_top_features.png}}
\caption{Mean absolute SHAP importance by source feature.}
\end{figure}

SHAP highlights months since previous record, tenure, and age as dominant descriptive signals. Compensation, contract type, illness, workload, and role-history features remain visible below those broad timing and life-cycle effects.

\begin{figure}[H]
\centering
\includegraphics[width=0.95\linewidth]{\figpath{permutation_reliance_file3.png}}
\caption{File3 permutation reliance using within-period shuffling.}
\end{figure}

The strongest external reliance is on tenure, age, and average payment. A positive AUC drop means performance decreased when that feature group was shuffled. Negative or near-zero values do not necessarily mean the feature is harmful; they may reflect noise, limited signal, or feature substitution by correlated variables.

\begin{figure}[H]
\centering
\includegraphics[width=0.95\linewidth]{\figpath{shap_beeswarm.png}}
\caption{SHAP beeswarm for the selected XGBoost model.}
\end{figure}

Each point is a scored record; horizontal position indicates whether the feature pushed risk upward or downward. Dense overlap and encoded categorical features make the figure best suited for pattern recognition rather than exact measurement.

\subsection{Domain and actionability summary}

Feature importance was grouped into business domains to avoid over-interpreting individual correlated variables. Several payment-history features describe overlapping aspects of compensation trajectory, and several illness features describe overlapping absence patterns.

\begin{table}[H]
\centering
\caption{SHAP importance by domain.}
\begin{tabular}{p{5cm}p{3.2cm}rr}
\toprule
Domain & Actionability & Mean abs SHAP & Share \\
\midrule
Data maturity / timing & Descriptive & 0.898 & 22.8\% \\
Compensation & Actionable & 0.687 & 17.5\% \\
Age / tenure & Descriptive & 0.654 & 16.6\% \\
Absence / wellness & Partially actionable & 0.536 & 13.6\% \\
Workload & Actionable & 0.535 & 13.6\% \\
Role / manager / organization & Partially actionable & 0.496 & 12.6\% \\
Location / commute & Descriptive & 0.081 & 2.1\% \\
Data availability & Descriptive & 0.045 & 1.1\% \\
\bottomrule
\end{tabular}
\end{table}

\begin{table}[H]
\centering
\caption{SHAP importance by actionability.}
\begin{tabular}{lrr}
\toprule
Actionability & Mean abs SHAP & Share \\
\midrule
Descriptive & 1.678 & 42.7\% \\
Actionable & 1.221 & 31.1\% \\
Partially actionable & 1.032 & 26.3\% \\
\bottomrule
\end{tabular}
\end{table}

\begin{figure}[H]
\centering
\includegraphics[width=0.95\linewidth]{\figpath{importance_by_domain.png}}
\caption{SHAP importance grouped by feature domain.}
\end{figure}

Data maturity/timing and age/tenure account for a large share of model behavior. Actionable areas remain meaningful: compensation, workload, illness/wellness, and role/manager/organization together represent a substantial share of model signal.

\begin{figure}[H]
\centering
\includegraphics[width=0.75\linewidth]{\figpath{importance_by_actionability.png}}
\caption{SHAP importance grouped by actionability.}
\end{figure}

Approximately 31\% of SHAP importance is from directly actionable features and about 26\% from partially actionable features. Descriptive variables are useful for segmentation but should not be treated as management levers.

\section{Actionable Importance Within Age and Tenure Groups}

Because age and tenure are strong descriptive predictors, a separate grouped analysis asked a more practical question: within comparable age or tenure groups, which actionable or partially actionable features still matter? This analysis used untouched file3 records, ranked actionable features by within-group SHAP, and measured within-group permutation PR-AUC drops when enough positives and negatives were available.

\subsection{Within age bands}

\begin{longtable}{lrrrlp{4cm}p{3cm}r}
\caption{Top actionable SHAP features within age bands.}\\
\toprule
Age band & Rows & Pos. & Prev. & Rank & Feature & Domain & SHAP \\
\midrule
\endfirsthead
\toprule
Age band & Rows & Pos. & Prev. & Rank & Feature & Domain & SHAP \\
\midrule
\endhead
18-34 & 721 & 76 & 10.5\% & 1 & \feature{avg\_Payment} & Compensation & 0.123 \\
18-34 & 721 & 76 & 10.5\% & 2 & \feature{contract\_type} & Role / manager / organization & 0.115 \\
18-34 & 721 & 76 & 10.5\% & 3 & \feature{tafkidCode\_hist\_prev} & Role / manager / organization & 0.111 \\
18-34 & 721 & 76 & 10.5\% & 4 & \feature{avg\_illness} & Absence / wellness & 0.089 \\
35-49 & 1,360 & 60 & 4.4\% & 1 & \feature{stdevp\_Payment\_hist\_prev} & Compensation & 0.109 \\
35-49 & 1,360 & 60 & 4.4\% & 2 & \feature{avg\_Payment} & Compensation & 0.107 \\
35-49 & 1,360 & 60 & 4.4\% & 3 & \feature{tafkidCode\_hist\_prev} & Role / manager / organization & 0.105 \\
35-49 & 1,360 & 60 & 4.4\% & 4 & \feature{contract\_type} & Role / manager / organization & 0.102 \\
50-59 & 710 & 14 & 2.0\% & 1 & \feature{stdevp\_Payment\_hist\_prev} & Compensation & 0.122 \\
50-59 & 710 & 14 & 2.0\% & 2 & \feature{tafkidCode\_hist\_prev} & Role / manager / organization & 0.108 \\
50-59 & 710 & 14 & 2.0\% & 3 & \feature{avg\_Payment} & Compensation & 0.100 \\
50-59 & 710 & 14 & 2.0\% & 4 & \feature{contract\_type} & Role / manager / organization & 0.083 \\
60+ & 349 & 3 & 0.9\% & 1 & \feature{tafkidCode\_hist\_prev} & Role / manager / organization & 0.159 \\
60+ & 349 & 3 & 0.9\% & 2 & \feature{contract\_type} & Role / manager / organization & 0.099 \\
60+ & 349 & 3 & 0.9\% & 3 & \feature{avg\_Payment} & Compensation & 0.099 \\
60+ & 349 & 3 & 0.9\% & 4 & \feature{illness\_cv} & Absence / wellness & 0.083 \\
\bottomrule
\end{longtable}

\begin{figure}[H]
\centering
\includegraphics[width=0.95\linewidth]{\figpath{actionable_shap_by_age_band.png}}
\caption{Top actionable SHAP features within age bands.}
\end{figure}

Average payment, payment variability, contract type, role history, illness, and workload remain visible after segmenting by age. The 60+ group has only three positives, so its ranking is descriptive and low-confidence.

\begin{table}[H]
\centering
\caption{Top actionable permutation PR-AUC drops within age bands.}
\begin{tabular}{llrrrr}
\toprule
Age band & Feature & Pos. & Prev. & Base PR-AUC & PR-AUC drop \\
\midrule
18-34 & \feature{avg\_Payment} & 76 & 10.5\% & 0.2135 & 0.0379 \\
18-34 & \feature{illness\_cv} & 76 & 10.5\% & 0.2135 & 0.0059 \\
18-34 & \feature{avg\_Payment\_hist\_prior\_std} & 76 & 10.5\% & 0.2135 & 0.0050 \\
35-49 & \feature{avg\_Payment} & 60 & 4.4\% & 0.0846 & 0.0087 \\
35-49 & \feature{stdevp\_Payment\_hist\_prev} & 60 & 4.4\% & 0.0846 & 0.0037 \\
35-49 & \feature{avg\_illness\_hist\_vs\_prior\_mean} & 60 & 4.4\% & 0.0846 & 0.0031 \\
50-59 & \feature{avg\_Payment} & 14 & 2.0\% & 0.0769 & 0.0157 \\
50-59 & \feature{illness\_cv} & 14 & 2.0\% & 0.0769 & 0.0124 \\
50-59 & \feature{avg\_Payment\_hist\_prior\_std} & 14 & 2.0\% & 0.0769 & 0.0035 \\
\bottomrule
\end{tabular}
\end{table}

\begin{figure}[H]
\centering
\includegraphics[width=0.95\linewidth]{\figpath{actionable_permutation_by_age_band.png}}
\caption{Within-age-band actionable permutation PR-AUC drops.}
\end{figure}

Average payment is the clearest actionable reliance signal in the 18-34, 35-49, and 50-59 age bands. The 60+ band is omitted from permutation interpretation because it had too few positives.

\subsection{Within tenure bands}

\begin{longtable}{lrrrlp{4cm}p{3cm}r}
\caption{Top actionable SHAP features within tenure bands.}\\
\toprule
Tenure & Rows & Pos. & Prev. & Rank & Feature & Domain & SHAP \\
\midrule
\endfirsthead
\toprule
Tenure & Rows & Pos. & Prev. & Rank & Feature & Domain & SHAP \\
\midrule
\endhead
$<2$y & 775 & 77 & 9.9\% & 1 & \feature{avg\_Payment} & Compensation & 0.142 \\
$<2$y & 775 & 77 & 9.9\% & 2 & \feature{contract\_type} & Role / manager / organization & 0.107 \\
$<2$y & 775 & 77 & 9.9\% & 3 & \feature{workload\_cv} & Workload & 0.094 \\
$<2$y & 775 & 77 & 9.9\% & 4 & \feature{avg\_illness} & Absence / wellness & 0.078 \\
2-5y & 749 & 39 & 5.2\% & 1 & \feature{avg\_Payment} & Compensation & 0.116 \\
2-5y & 749 & 39 & 5.2\% & 2 & \feature{stdevp\_Payment\_hist\_prev} & Compensation & 0.115 \\
2-5y & 749 & 39 & 5.2\% & 3 & \feature{tafkidCode\_hist\_prev} & Role / manager / organization & 0.115 \\
2-5y & 749 & 39 & 5.2\% & 4 & \feature{contract\_type} & Role / manager / organization & 0.108 \\
5-10y & 795 & 24 & 3.0\% & 1 & \feature{stdevp\_Payment\_hist\_prev} & Compensation & 0.146 \\
5-10y & 795 & 24 & 3.0\% & 2 & \feature{tafkidCode\_hist\_prev} & Role / manager / organization & 0.141 \\
5-10y & 795 & 24 & 3.0\% & 3 & \feature{avg\_Payment\_hist\_prior\_std} & Compensation & 0.096 \\
5-10y & 795 & 24 & 3.0\% & 4 & \feature{contract\_type} & Role / manager / organization & 0.094 \\
10y+ & 821 & 13 & 1.6\% & 1 & \feature{stdevp\_Payment\_hist\_prev} & Compensation & 0.125 \\
10y+ & 821 & 13 & 1.6\% & 2 & \feature{tafkidCode\_hist\_prev} & Role / manager / organization & 0.123 \\
10y+ & 821 & 13 & 1.6\% & 3 & \feature{contract\_type} & Role / manager / organization & 0.094 \\
10y+ & 821 & 13 & 1.6\% & 4 & \feature{avg\_Payment\_hist\_delta\_prev} & Compensation & 0.092 \\
\bottomrule
\end{longtable}

\begin{figure}[H]
\centering
\includegraphics[width=0.95\linewidth]{\figpath{actionable_shap_by_tenure_band.png}}
\caption{Top actionable SHAP features within tenure bands.}
\end{figure}

For employees with less than two years tenure, average payment is the leading actionable feature. For 5-10 and 10+ years tenure, payment variability/history, role history, illness history, and contract type become more prominent.

\begin{table}[H]
\centering
\caption{Top actionable permutation PR-AUC drops within tenure bands.}
\begin{tabular}{llrrrr}
\toprule
Tenure & Feature & Pos. & Prev. & Base PR-AUC & PR-AUC drop \\
\midrule
$<2$y & \feature{avg\_Payment} & 77 & 9.9\% & 0.2294 & 0.0488 \\
$<2$y & \feature{illness\_cv} & 77 & 9.9\% & 0.2294 & 0.0029 \\
$<2$y & \feature{contract\_type} & 77 & 9.9\% & 0.2294 & 0.0028 \\
2-5y & \feature{illness\_cv} & 39 & 5.2\% & 0.0761 & 0.0052 \\
2-5y & \feature{avg\_Payment} & 39 & 5.2\% & 0.0761 & 0.0036 \\
2-5y & \feature{avg\_Payment\_hist\_prior\_std} & 39 & 5.2\% & 0.0761 & 0.0027 \\
5-10y & \feature{avg\_illness} & 24 & 3.0\% & 0.1418 & 0.0264 \\
5-10y & \feature{avg\_Payment} & 24 & 3.0\% & 0.1418 & 0.0247 \\
5-10y & \feature{avg\_Payment\_hist\_delta\_prev} & 24 & 3.0\% & 0.1418 & 0.0063 \\
10y+ & \feature{avg\_illness\_hist\_vs\_prior\_mean} & 13 & 1.6\% & 0.0759 & 0.0053 \\
10y+ & \feature{stdevp\_Payment\_hist\_prev} & 13 & 1.6\% & 0.0759 & 0.0040 \\
10y+ & \feature{tafkidCode\_hist\_prev} & 13 & 1.6\% & 0.0759 & 0.0039 \\
\bottomrule
\end{tabular}
\end{table}

\begin{figure}[H]
\centering
\includegraphics[width=0.95\linewidth]{\figpath{actionable_permutation_by_tenure_band.png}}
\caption{Within-tenure-band actionable permutation PR-AUC drops.}
\end{figure}

The permutation results suggest different levers by tenure stage: payment dominates among newer employees; illness and payment both matter in the 5-10 year band; and illness-history and payment-history features matter most for long-tenure employees. Some negative or near-zero drops in the full CSV should be treated as weak evidence rather than evidence of causal protection.

\section{Streamlit Dashboard Integration}

The final model was incorporated into the existing Streamlit executive dashboard while preserving the original visual style and tab structure. A new final-model dataset option builds dashboard-compatible latest-employee records from final file3 predictions. The micro-level what-if simulator was revised so that it reflects the final feature-importance results rather than merely showing percent input changes.

Dashboard changes:
\begin{itemize}
    \item Macro and meso views continue to summarize turnover probability, risk categories, budget-section risk, tenure-risk patterns, and highest-risk employees.
    \item The micro view now supports the final artifact and scores prepared final-model rows using the locked XGBoost pipeline.
    \item What-if levers focus on salary/payment, workload, illness, contract type, job rank, and budget section.
    \item For final-model mode, the lever-impact chart reports estimated risk-point impact from each lever, aligning the interface more closely with model interpretability outputs.
\end{itemize}

\section{Limitations and Recommended Next Steps}

The model demonstrates useful ranking performance, but several limitations remain. These are methodological constraints imposed by the data and prediction problem rather than simply implementation defects.

\begin{itemize}
    \item \textbf{Generalizability:} File3 differs from file1/file2 in prevalence and likely feature distributions. The model may learn source-specific patterns that transfer imperfectly.
    \item \textbf{Payment imputation:} File2 payment imputation is noisy. Payment is important to the final model, so systematic imputation error can affect both performance and explanations.
    \item \textbf{Low event prevalence:} File3 prevalence is only 4.87\%, so precision remains modest even when ranking metrics are above baseline.
    \item \textbf{Small subgroups:} The 60+ age group contains only three positives in file3, making subgroup permutation PR-AUC unreliable.
    \item \textbf{Causality:} Feature importance is predictive, not causal. Compensation, workload, and illness features can guide review, but interventions should be validated with domain knowledge and policy constraints.
    \item \textbf{Temporal framing:} A future improvement would define a prediction horizon, such as resignation within the next 3, 6, or 12 months, rather than treating every row as an undifferentiated turnover prediction row.
\end{itemize}

\textbf{Recommended next research step.} The strongest next methodological improvement is to quantify source drift between file1/file2 and file3 and then test portable/actionable feature sets under source-aware validation. This would clarify whether weaker PR-AUC is caused by class imbalance, missingness, population shift, or unstable company-specific codes.

Additional useful extensions:
\begin{enumerate}
    \item Reframe the target as resignation within a fixed future horizon, such as 3, 6, or 12 months.
    \item Compare row-level predictions against latest-row-per-employee predictions, since HR interventions usually occur at the employee level.
    \item Run train-file1/validate-file2 and train-file2/validate-file1 experiments to quantify cross-source portability.
    \item Compare feature sets that remove company-specific codes, age, tenure, and location to determine the tradeoff between predictive performance and generalizability.
    \item Evaluate calibration on file3, especially if probabilities are used directly in the dashboard rather than only for ranking.
\end{enumerate}

\appendix
\section{Files Produced by the Final Pipeline}
\begin{itemize}
    \item Model artifact: \feature{artifacts/final\_best\_model.pkl}
    \item Predictions: \feature{output/final/file3\_predictions.xlsx}
    \item Model comparison: \feature{output/final/model\_comparison.csv}
    \item Feature-importance report: \feature{output/final/feature\_importance/feature\_importance\_report.txt}
    \item Grouped actionable report: \feature{output/final/feature\_importance/grouped\_actionable/grouped\_actionable\_importance\_report.txt}
    \item Streamlit integration helpers: \feature{src/final\_dashboard.py} and updated \feature{src/inference.py}
\end{itemize}

\section{Metric Definitions}
\begin{itemize}
    \item \textbf{ROC-AUC:} Measures how often the model ranks a positive example above a negative example. It is useful for ranking but can look optimistic under heavy class imbalance.
    \item \textbf{PR-AUC:} Summarizes precision-recall behavior and is more sensitive to rare-event prediction quality. It is the primary selection metric in this final pipeline.
    \item \textbf{F1:} Harmonic mean of precision and recall at a selected threshold.
    \item \textbf{Recall@Top20\%:} Share of positives captured when reviewing the top 20\% highest-risk rows.
    \item \textbf{Precision@Top20\%:} Positive rate among the top 20\% highest-risk rows.
    \item \textbf{Brier score:} Probability calibration error; lower is better.
    \item \textbf{Log loss:} Penalizes confident wrong probabilities; lower is better.
\end{itemize}

\end{document}


